%%%%%%%%%%  scan idx 644  |  folio 631  %%%%%%%%%%
\SGAPrintPage{631}%
we have: \(a_i^{(m)}=P_i^{(m,r)}(a_r,\ldots,a_i)\).

\sgaStableTarget{sga6.e13.lem.2-11}\textbf{Lemma 2.11.} \emph{Let \(S\) be a noetherian scheme and \(X\) a projective \(S\)-scheme, endowed with a very ample sheaf \(H={O}_X(1)\). We suppose that the fibers of \(X/S\) are geometrically integral, of the same dimension \(r\). Let \(\mathcal L\) be a family of classes of invertible sheaves \(L_K\) on the fibers of \(X/S\). Then for \(\mathcal L\) to be bounded, it is (necessary and) sufficient that, in the Hilbert polynomials \(\chi(L_K(n))=\sum_{i=0}^{r}a_i\binom{n+i}{i}\) of the \(L_K\in\mathcal L\), all the coefficients \(a_i\) remain bounded. Moreover, if the fibers of \(X/S\) are also geometrically normal, then it suffices that the three coefficients \(a_r,a_{r-1},a_{r-2}\) remain bounded.}

Indeed, let \(\delta=\sup\{d(L_K)\}\). If \(\delta>0\), one may replace \(\mathcal L\) by the family of classes of the \(L_K(-\delta)\) by 2.3 and 2.10\,; one may thus suppose \(\delta\le0\). Then \(\mathcal L\) is a \((b)\)-family by 2.5, hence bounded if all the coefficients \(a_i\) remain bounded, by 1.13, and, when the fibers of \(X/S\) are geometrically normal, it suffices for this that the three coefficients \(a_r,a_{r-1},a_{r-2}\) remain bounded, by 2.9.

\sgaStableTarget{sga6.e13.sec.3}\subsection*{3. General finiteness theorems}
\sgaStableTarget{sga6.e13.lem.3-1}\textbf{Lemma 3.1.} \emph{Let \(S\) be a scheme and \(X\) a quasi-compact and quasi-separated \(S\)-scheme\,; let \(\underline{\Pic}_{X/S}\) be the Picard functor, (localized for the f.p.p.f. topology)}\SGAAuthorFootnote{*}{Cf. SGA 3 IV 6.3, and FGA no.\ 232 for the definition of \(\underline{\Pic}_{X/S}\), where, as indicated, the fpqc topology should here be replaced by the fppf topology.}.
\emph{Then for every filtered inverse system \((Z_\alpha)\) of quasi-compact and quasi-separated \(S\)-schemes
%%%%%%%%%%  scan idx 645  |  folio 632  %%%%%%%%%%
\SGAPrintPage{632}%
with affine transition morphisms, the canonical map}
\sgaStableTarget{sga6.e13.eq.3-1-1}\[
\varinjlim \underline{\Pic}_{X/S}(Z_\alpha)\longrightarrow
\underline{\Pic}_{X/S}\bigl(\varprojlim Z_\alpha\bigr)
\tag{3.1.1}
\]
\emph{is bijective.}

\emph{Consequently} (EGA IV 8.14.2), \emph{if the Picard functor is represented by an \(S\)-scheme \(\Pic_{X/S}\), then \(\Pic_{X/S}\) is locally of finite presentation.}

Indeed, an element \(\lambda_\alpha\) of \(\underline{\Pic}_{X/S}(Z_\alpha)\) is represented by an invertible sheaf \(L'_\alpha\) on an \(X_{Z'_\alpha}\) where \(Z'_\alpha\to Z_\alpha\) is an f.p.p.f. morphism. Put \(Z=\varprojlim Z_\alpha\). The image \(\lambda\) of \(\lambda_\alpha\) in \(\underline{\Pic}_{X/S}(Z)\) is represented by the invertible sheaf \(L'\) on \(X_{Z'}\), the pullback of \(L'_\alpha\), where \(Z'=Z'_\alpha\times_{Z_\alpha}Z\).

If \(\lambda\) is the identity element of \(\underline{\Pic}_{X/S}(Z)\), there exists an f.p.p.f. morphism \(Z''\to Z'\) such that the pullback of \(L'\) to \(X_{Z''}\) is \(\simeq{O}_{X_{Z''}}\). By (EGA IV 8.8.2, 8.10.5 (vi), 11.2.6 (ii), 8.5.2 (i), and 8.5.2.4), there exist a \(\beta\geq\alpha\) and an f.p.p.f. morphism \(Z''_\beta\to Z'_\beta\), where \(Z'_\beta=Z'_\alpha\times_{Z_\alpha}Z_\beta\), such that the pullback of \(L'_\alpha\) to \(X_{Z''_\beta}\) is \(\simeq{O}_{X_{Z''_\beta}}\). The image of \(\lambda_\alpha\) is therefore already the identity element in \(\underline{\Pic}_{X/S}(Z_\beta)\). Consequently, 3.1.1 is injective.

The surjectivity of 3.1.1 is proved in an analogous way.

\sgaStableTarget{sga6.e13.prop.3-2}\textbf{Proposition 3.2.} \emph{Let \(S\) be a noetherian scheme and \(X\) a proper \(S\)-scheme. We suppose that the Picard functor is represented by an \(S\)-scheme \(\Pic_{X/S}\). Then}~:
\begin{enumerate}[label=(\roman*)]
\item \emph{\(\Pic_{X/S}\) is locally of finite type over \(S\).}
%%%%%%%%%%  scan idx 646  |  folio 633  %%%%%%%%%%
\item\SGAPrintPage{633}%
\emph{The subsets \(\Lambda\) of \(\Pic_{X/S}\) correspond bijectively to the families \(\mathscr L\) of classes of invertible sheaves on the fibers of \(X/S\).}
\item \emph{A family \(\Lambda\) of \(\Pic_{X/S}\) is quasi-compact if and only if the corresponding family \(\mathscr L\) is bounded.}
\end{enumerate}

Indeed, (i) follows from (3.1)\,; (ii), from the following fact: if \(s\) is a point of \(S\) and \(K\) is an algebraically closed extension of \(k(s)\), then the \(K\)-points of \(\Pic_{X/S}\) correspond bijectively to the isomorphism classes of invertible sheaves on \(X_K\).

As for (iii), suppose that the subset \(\Lambda\) of \(\Pic_{X/S}\) is quasi-compact. Then \(\Lambda\) is contained in an open subset \(U\) of finite type, and the inclusion \(U\to\Pic_{X/S}\) corresponds to an invertible sheaf \(L\) on an \(X_T\), where \(T\to U\) is a surjective, (flat) morphism of finite type. Finally, \(T\to S\) is of finite type, so \(L\) bounds \(\mathscr L\).

Conversely, suppose that \(\mathscr L\) is bounded by a coherent sheaf \(L\) on an \(X_T\), with \(T\) of finite type over \(S\). One may suppose \(L\) invertible (1.13 (i)) so that \(L\) defines a morphism \(f:T\to\Pic_{X/S}\). The image of \(f\) contains \(\Lambda\) and is quasi-compact, hence noetherian by (i). Consequently, \(\Lambda\) is quasi-compact, and we have (iii).

By 3.2, the usual definition of morphisms of finite type between two Picard schemes generalizes to the following:

\sgaStableTarget{sga6.e13.def.3-3}\textbf{Definition 3.3.} \emph{Let \(X\) and \(Y\) be proper over noetherian \(S\). A morphism between the Picard functors}
\[
\begin{tikzcd}
\underline{\Pic}_{Y/S}\arrow[r]&\underline{\Pic}_{X/S}
\end{tikzcd}
\]
%%%%%%%%%%  scan idx 647  |  folio 634  %%%%%%%%%%
\SGAPrintPage{634}%
\emph{is said to be} of finite type \emph{if every bounded family of classes of invertible sheaves on the fibers of \(X/S\) has as its inverse image a bounded family of classes of sheaves on the fibers of \(Y/S\).}

\sgaStableTarget{sga6.e13.lem.3-4}\textbf{Lemma 3.4.} \emph{Let \(X,Y,X',Y'\) be proper over noetherian \(S\), and \(T\to S\) surjective of finite type. Then}~:
\begin{enumerate}[label=(\roman*)]
\item \emph{For a morphism \(\Phi:\underline{\Pic}_{Y/S}\to\underline{\Pic}_{X/S}\) to be of finite type, it is necessary and sufficient that \(\Phi_T:\underline{\Pic}_{Y_T/T}\to\underline{\Pic}_{X_T/T}\) be of finite type.}
\item \emph{In a commutative diagram}
\[
\begin{tikzcd}
\Phi':\underline{\Pic}_{Y'/S}\arrow[r]&
\underline{\Pic}_{X'/S}\\
\Phi:\underline{\Pic}_{Y/S}\arrow[r]\arrow[u,"\Phi_Y"]&
\underline{\Pic}_{X/S}\arrow[u,"\Phi_X"]
\end{tikzcd}
\]
\emph{with \(\Phi'\) and \(\Phi_Y\) of finite type, \(\Phi\) is of finite type.}
\end{enumerate}

Indeed, suppose that \(\Phi_T\) is of finite type. Let \(\mathscr L\) be a family on the fibers of \(X/S\), bounded by an \(L\) on an \(X_R\). The family \(\mathscr L_T=\mathscr L\times_S T\) on the fibers of \(X_T/T\), defined in the evident way, is bounded by the pullback of \(L\) to \(X_{R\times T}\). If \(\Phi_T^{-1}(\mathscr L_T)\) is bounded by an \(M\) on a \(Y_Q\), then \(M\) also bounds \(\Phi^{-1}(\mathscr L)\).

The converse of (i) and assertion (ii) are proved in an analogous way.

\sgaStableTarget{sga6.e13.thm.3-5}\textbf{Theorem 3.5.} \emph{Let \(S\) be a noetherian scheme, \(X\) and \(Y\) two proper \(S\)-schemes and \(f:X\to Y\) a morphism. If \(f\) is surjective, then the induced morphism between the Picard functors}
\[
f^*:\underline{\Pic}_{Y/S}\longrightarrow\underline{\Pic}_{X/S}
\]
%%%%%%%%%%  scan idx 648  |  folio 635  %%%%%%%%%%
\SGAPrintPage{635}%
\emph{is of finite type.}

Indeed, by 3.4 (i), one may evidently suppose \(S\) reduced and irreducible. Then there exists a nonempty open subset \(U\) of \(S\) such that \(\underline{\Pic}_{Y/S}\) and \(\underline{\Pic}_{X/S}\) are representable and the restriction of \(f^*\) is quasi-affine (XII 1.1), hence of finite type by 3.2 (i). Endowing \(T=S-U\) with the induced reduced structure, one may suppose that the morphism \(\underline{\Pic}_{Y_T/T}\longrightarrow\underline{\Pic}_{X_T/T}\) induced by \(f^*\) is of finite type by noetherian induction\,; this gives the assertion, by 3.4 (i) applied to \(U\amalg T\to S\).

\sgaStableTarget{sga6.e13.thm.3-6}\textbf{Theorem 3.6.} \emph{Let \(S\) be a noetherian scheme, \(X\) a proper \(S\)-scheme and \(\underline{\Pic}_{X/S}\) the Picard functor. Then, for every integer \(n\ne0\), the \(n\)-th power homomorphism}
\[
\Phi_n:\underline{\Pic}_{X/S}\longrightarrow\underline{\Pic}_{X/S}
\]
\emph{is a morphism of finite type.}

Indeed, let \(f:X'\to X\) be a Chow presentation: \(X'\) is projective over \(S\) and \(f\) is surjective. If the theorem is admitted for \(X'\), then the theorem for \(X\) follows by 3.5 and 3.4 (ii). One may therefore suppose \(X/S\) projective.

There exist a nonempty open subset \(U\) of \(S\), an integral scheme \(S'\) and a finite surjective morphism \(S'\to U\), such that every irreducible component \(X'_i\) of \(X_{S'}\), endowed with the induced reduced structure, is integral over \(S'\) (XII 1.3). Endowing \(T=S-U\) with the induced reduced structure, one may admit the theorem for \(X_T/T\) by noetherian induction. If it is admitted for all the \(X'_i/S'\), it follows for \(\coprod X'_i/S'\) by the
